\documentclass[11pt]{article}
\usepackage[a4paper,margin=25mm]{geometry}
\usepackage{amsmath,amssymb,amsthm,mathtools,microtype,needspace}
\usepackage[colorlinks=true,linkcolor=blue,citecolor=blue,urlcolor=blue]{hyperref}
\newtheorem{theorem}{Theorem}
\newtheorem{lemma}[theorem]{Lemma}
\newtheorem{corollary}[theorem]{Corollary}
\newtheorem{proposition}[theorem]{Proposition}
\theoremstyle{remark}\newtheorem{remark}[theorem]{Remark}
\newcommand{\D}{\mathbb D}
\newcommand{\C}{\mathbb C}
\newcommand{\norm}[1]{\lVert #1\rVert}
\newcommand{\Rea}{\operatorname{Re}}
\newcommand{\ee}{\mathrm e}
\title{A sharp circle-root residual barrier\\and its collapse with one interior zero}
\author{Lluis Eriksson}
\date{18 October 2026}
\begin{document}
\maketitle
\begin{abstract}
For degree-\(n\) polynomials with all zeros on the unit circle,
the least normalized residual \(\norm{P'-\lambda P}_{H^2}/\norm P_{H^2}\)
tends to \(1/\sqrt8\), both for arbitrary complex \(\lambda\) and
for the restricted slope \(\lambda=P'(0)\), \(P(0)=1\).
We prove the sharp lower limit through the classical Mathieu ground
state and give an explicit circle-root family attaining it in every
large degree. In contrast, exactly one zero strictly inside the
disk, with all remaining zeros on the circle, permits restricted
residual \(O(n^{-1/6})\), giving the exact zero-count threshold
for collapse. More generally, for each fixed \(A>0\),
a fixed number \(M_A\) of interior zeros permits \(O(n^{-A})\).
Their distance and radial gap from 1 have order \(n^{-1/3}\);
explicit algebraic parameters in the one-root case, and an
implicit-function construction in general, correct both degree rounding
and the exact slope. For all disk-root polynomials and arbitrary slope,
we also prove a uniform lower bound \(\exp(-C\sqrt n)\) for every
\(C>\sqrt{(8/3)\log17}\), explicitly \(\exp(-3\sqrt n)\) for
\(n\ge10^4\). This refines the known residual method for exterior
power sums; the negative answer to Erd\H{o}s Problem 973 was already
known. The optimal disk and one-root decay rates, and small exterior
power-sum constructions, remain unresolved.
\end{abstract}

\section{Statement and relation to previous work}
All logarithms are natural. The measure on the unit circle is normalized
to have mass one. Thus for \(P(z)=\sum_{j=0}^na_jz^j\),
\[
 \norm P_2^2=\frac1{2\pi}\int_0^{2\pi}|P(\ee^{it})|^2\,dt
             =\sum_{j=0}^n|a_j|^2,
 \qquad \norm P_\infty=\max_{|z|=1}|P(z)|.
\]
Write \(\mathcal S_n\) for the nonzero polynomials of exact degree \(n\)
whose zeros are in \(\overline\D\), including multiplicity and boundary
zeros. Define
\[
 s_n=\inf_{\substack{P\in\mathcal S_n\\\lambda\in\C}}
          \frac{\norm{P'-\lambda P}_2}{\norm P_2}.
\]
There is no normalization at the origin in this definition.

\begin{theorem}\label{thm:main}
Put \(C_0=\sqrt{(8/3)\log17}\). For every \(C>C_0\) there is \(N_C\)
such that
\begin{equation}\label{eq:asymptotic}
 s_n\ge \ee^{-C\sqrt n}\qquad(n\ge N_C).
\end{equation}
The following explicit bound also holds:
\begin{equation}\label{eq:explicit}
 s_n\ge \ee^{-3\sqrt n}\qquad(n\ge10^4).
\end{equation}
\end{theorem}

For comparison, Tan, Wang, Huang and Chen \cite[Theorem 1.1]{tan}
considered the restricted infimum
\[
 r_n=\inf_{\substack{P\in\mathcal S_n\\P(0)=1}}
       \frac{\norm{P'-P'(0)P}_2}{\norm P_2}
\]
and proved \(r_n\ge\exp(-(1+o(1))\sqrt n\log n)\).
For an individual normalized polynomial we write
\(\rho(P)=\norm{P'-P'(0)P}_2/\norm P_2\).
Theorem~\ref{thm:circle} identifies the positive limiting residual
when every root lies on the circle. Theorem~\ref{thm:one} shows
that one interior root suffices to make it vanish, establishing the
minimal zero-count threshold. These statements
concern the global weighted boundary norm, including the exact
restricted slope.

Since \(s_n\le r_n\), Theorem~\ref{thm:main} strengthens that quantitative
estimate and allows an arbitrary scalar slope.
Their maximum-point normalization, segment integration and projected
Newton identity are retained here, with attribution. The two changes
are a local analytic bound for the logarithmic derivative and a bounded
test polynomial. These remove the logarithmic factor in the exponent.

Luo, Yang and Zhu \cite[Theorem 1.1]{luo} had already proved that exterior
power sums cannot be exponentially small in the degree, answering
Erd\H{o}s Problem 973 negatively. The application in Section~\ref{sec:power}
is a quantitative refinement of that established conclusion.
No new solution of the original yes-or-no problem is asserted.

\begin{theorem}[Superalgebraic decay]\label{thm:upper}
For every \(A>0\), there are constants \(C_A,N_A\) such that
\[
 r_n\le C_A n^{-A}\qquad\text{for every integer }n\ge N_A.
\]
The polynomials can be chosen with real coefficients, exactly
\(M_A\) zeros strictly inside the unit disk, where \(M_A\) is
independent of \(n\), and all other zeros on its boundary.
The interior zeros approach 1 with distance and radial gap
\(\Theta(n^{-1/3})\). Together with
Theorem~\ref{thm:main}, this implies superalgebraic decay of both
\(s_n,r_n\), while \(s_n^{1/n},r_n^{1/n}\to1\).
\end{theorem}

The proof in Section~\ref{sec:upper} fixes the cancellation order
and finite integer multiplicities first, then lets the reciprocal root
scale grow as \(n^{1/3}\). It provides no effective bounds on
\(M_A,C_A,N_A\) or uniform control as \(A\) grows.
Section~\ref{sec:radial} explains why a fixed strict root radius
precludes superalgebraic decay of a single sequence.

Approximation by logarithmic derivatives and finite moment realization
have substantial prior history. Komarov \cite{komarovrate} obtains
geometric approximation on each fixed subdisk by logarithmic
derivatives of circle-root polynomials; \cite{komarovmoments}
also realizes prescribed initial moments with unit weights.
Those interior estimates do not supply the global weighted boundary
norm used here, with a target slope growing with degree. For example,
\(1+z^n\) approximates zero logarithmically on subdisks but has
restricted residual \(n/\sqrt2\). The present contribution is the
combination of exact slope, global concentration and correction at
every large degree, rather than moment realization alone.


\section{The sharp barrier when every zero is on the circle}\label{sec:circle}
Let \(s_n^\circ\) and \(r_n^\circ\) denote the same infima as
\(s_n,r_n\), with every zero required to lie on \(|z|=1\).

\begin{theorem}\label{thm:circle}
The unrestricted and restricted circle-root residuals satisfy
\[
 \lim_{n\to\infty}s_n^\circ
 =\lim_{n\to\infty}r_n^\circ=\frac1{\sqrt8}.
\]
The lower asymptotic bound also holds in the larger class of
self-inversive polynomials. No finite-degree extremizer is asserted.
\end{theorem}

\begin{proof}
We first prove the lower bound, allowing an arbitrary complex slope.
Rotating the variable makes this slope real and nonnegative, denoted
\(\lambda\). Self-inversivity permits a constant phase choice such that
\[
 P(\ee^{it})=\ee^{int/2}p(t),\qquad p(t)\in\mathbb R,
 \qquad p(t+2\pi)=(-1)^n p(t).
\]
These operations preserve the unrestricted normalized residual.
The familiar phase derivative \(n/2\) for circle-root polynomials
appears explicitly in \cite[(1.3)--(1.4)]{erdelyi}; here we retain
the additional scalar subtraction. Direct differentiation gives
\[
 P'-\lambda P=\ee^{i(n/2-1)t}
 \bigl[(n/2-\lambda\cos t)p-i(p'+\lambda\sin t\,p)\bigr].
\]
Integration of the cross term by parts is legitimate because
\(p^2\) is periodic. Consequently the squared residual is exactly
\begin{equation}\label{eq:circle-energy}
 \frac{\int_0^{2\pi}\!
 [p'^2+(n^2/4+\lambda^2-(n+1)\lambda\cos t)p^2]\,dt}
 {\int_0^{2\pi}p^2\,dt}.
\end{equation}
Write \(E(L)\) for the lowest periodic eigenvalue of
\(-d^2/dt^2-L\cos t\), for \(L\ge0\). In odd degree, \(|p|\)
is periodic in \(H^1\) and has the same quadratic-form values as
\(p\). Thus in both parities \eqref{eq:circle-energy} is at least
\[
 g_n(\lambda)=n^2/4+\lambda^2+E((n+1)\lambda).
\]
This is a variational lower bound, not an identification of the
periodic and antiperiodic spectra.

The change of variables \(t=2x+\pi\) gives
\(E(L)=a_0(2L)/4\), in the standard Mathieu convention
\(-d^2/dx^2+2q\cos(2x)\); see \cite[\S28.2]{dlmf}.
The use of this ground state in circular variational problems is
classical; compare \cite[\S VI.A.2]{busch}. The fixed-index expansion
\cite[(28.8.1), \(m=0,s=1\)]{dlmf} is
\[
 a_0(q)=-2q+2\sqrt q-\tfrac14+O(q^{-1/2}),
 \qquad
 E(L)=-L+\sqrt{L/2}-\tfrac1{16}+O(L^{-1/2}).
\]
We must minimize uniformly over \(\lambda\), rather than applying
this expansion only at a fixed shift. The variational principle
gives \(E(L)\ge-L\) and Lipschitz continuity, so \(g_n\) is
continuous and coercive, with
\[
 g_n(\lambda)\ge
 (\lambda-(n+1)/2)^2-(2n+1)/4.
\]
The trial value \(\lambda=n/2+1/4\) gives
\(g_n=1/8+O(n^{-1})\). Every minimizer therefore has the form
\(n/2+c_n\), where \(c_n=O(\sqrt n)\). On this window
\((n+1)\lambda\asymp n^2\), and the displayed single-variable
Mathieu remainder is uniformly \(O(n^{-1})\). Expanding its square
root uniformly on the same window yields
\[
 g_n(n/2+c)=(c-1/4)^2+\tfrac18+O((c^2+1)/n).
\]
The upper bound on the minimum and absorption of the \(c_n^2/n\)
term imply \(c_n=O(1)\). Applying the expansion again proves
\(\inf_{\lambda\ge0}g_n(\lambda)=1/8+O(n^{-1})\).
This supplies the required lower limit, also for the restricted slope.

For the upper bound, write \(n=6m+k\), \(0\le k\le5\), and for
\(m\ge3\) choose
\begin{align*}
 c_{m,k}=\cos t_{m,k}&=\tfrac12+\frac{2k-1}{8m},\\
 B_t(z)&=(1+z)^4(1-2\cos t\,z+z^2),\\
 P_{m,k}(z)&=B_{t_{m,k}}(z)^m(1+z)^k.
\end{align*}
All zeros lie on the circle, \(P_{m,k}(0)=1\), and
\[
 P_{m,k}'(0)=n/2+1/4,\qquad
 P_{m,k}'(1)/P_{m,k}(1)=n/2.
\]
The limiting factor is \(B_*=(1+z)^4(1-z+z^2)\).
With \(s=\sin^2(\theta/2)\), its squared normalized modulus is
\(\cos^6(\theta/2)\cos^2(3\theta/2)\), with a unique maximum at
\(\theta=0\), and its logarithm is \(-3\theta^2+O(\theta^4)\).
More explicitly, for the perturbed factor this modulus is
\[
 H_c(s)=(1-s)^4\left(1-\frac{2s}{1-c}\right)^2.
\]
Its secondary maximum has value
\((1+c)^6/[729(1-c)^2]<1\) for \(0\le c\le3/4\).
All six parameter sequences eventually lie in this interval.
Local negative curvature and the strict outer gap are uniform as
\(c\to1/2\). They give Gaussian domination near zero and an
exponentially small outer tail.

For \(F_{m,k}=P_{m,k}'/P_{m,k}-P_{m,k}'(0)\), the identity
\[
 B_*'/B_*-3=-\frac{3z(z-1)^2}{z^3+1}
\]
shows, on putting \(\theta=u/\sqrt m\), that
\[
 F_{m,k}(\ee^{i\theta})\longrightarrow-\tfrac14+\tfrac32u^2,
 \qquad
 |P_{m,k}(\ee^{i\theta})/P_{m,k}(1)|^2\longrightarrow\ee^{-3u^2}.
\]
Indeed the constant term of \(F_{m,k}\) at 1 is exactly \(-1/4\),
its linear term is \(O(1)\), and its quadratic term is
\(-3m/2+O(1)\). A fixed zero-free neighborhood gives the local bound
\(C(1+|\theta|+m\theta^2)\); after rescaling, a polynomial times
a Gaussian dominates. On the outer interval, differentiate the
polynomial product, removing one factor. Uniformly bounded normalized
factor derivatives give an outer squared numerator
\(O(m^2q^{2m-2})\), for some \(q<1\). Thus no estimate at a
pole of the logarithmic derivative is required.
The ratio of the complete integrals tends to
\[
 \frac{\int_{\mathbb R}(-1/4+3u^2/2)^2\ee^{-3u^2}\,du}
 {\int_{\mathbb R}\ee^{-3u^2}\,du}
 =\tfrac1{16}-\tfrac34\cdot\tfrac16+\tfrac94\cdot\tfrac1{12}
 =\tfrac18.
\]
Uniformity over the six values of \(k\) covers every sufficiently
large degree. This proves the upper bound for \(r_n^\circ\), and
\(s_n^\circ\le r_n^\circ\) completes both limits.
\end{proof}


\section{A finite residual certificate}
We first give conditions that can be checked at a fixed degree.
The asymptotic and explicit choices will then be elementary.

\begin{lemma}[Normalization and separation]\label{lem:normalize}
Suppose \(P\in\mathcal S_n\), \(\lambda\in\C\), \(y>0\), and
\[
 \frac{\norm{P'-\lambda P}_2}{\norm P_2}\le\ee^{-y},
 \qquad \sqrt{n+1}\,\ee^{-y}\le1.
\]
Set
\begin{equation}\label{eq:eta}
 \gamma_n=2+\log(2\sqrt{n+1}),\qquad
 \eta=\frac{y-\gamma_n}{n}.
\end{equation}
If \(\eta>0\), there are numbers \(v_1,\ldots,v_n\) and moments
\(m_k=n^{-1}\sum_jv_j^k\) such that
\begin{equation}\label{eq:geometry}
 m_1\in\mathbb R,\quad m_1\ge\tfrac12,\qquad
 \Rea v_j\ge\tfrac12,\quad |v_j|^2\le\frac{\Rea v_j}{\eta}.
\end{equation}
They are the reciprocals \(v_j=(1-\beta_j)^{-1}\) of the complex displacements
from the zeros to a normalized maximum point.
\end{lemma}
\begin{proof}
This is the normalization and separation argument of
\cite[Sections 2.1--2.2]{tan}, with the scalar slope left arbitrary.
Choose \(|\tau|=1\) at a maximum of \(|P|\), and put
\[
 G(w)=\frac{P(\tau w)}{P(\tau)},\qquad
 \mu=\tau\lambda,\qquad E=G'-\mu G.
\]
The denominator is nonzero because \(P\) is nonzero. We have
\(G(1)=\norm G_\infty=1\), \(\norm G_2\le1\), and
\(\norm E_2\le\ee^{-y}\).
Differentiating \(|G(\ee^{it})|^2\) at its maximum shows that
\(d=G'(1)\) is real. Bernstein's polynomial inequality gives
\(|d|\le n\). The coefficient bound
\(\norm E_\infty\le\sqrt{n+1}\norm E_2\) therefore yields
\begin{equation}\label{eq:mu}
 |\mu-d|=|E(1)|\le1,\qquad |\mu|\le n+1.
\end{equation}

Let \(\beta\) be any zero of \(G\), and set \(\delta=1-\Rea\beta\).
Integrating the exact ODE along the segment from \(\beta\) to \(1\) gives
\[
 1=\int_\beta^1\ee^{\mu(1-w)}E(w)\,dw.
\]
That segment lies in \(\overline\D\). On it,
\(\Rea(\mu(1-w))\le n\delta+2\), using \(d\in\mathbb R\),
\(|d|\le n\), \(|\mu-d|\le1\), and \(|1-\beta|\le2\).
Consequently
\[
 1\le 2\sqrt{n+1}\,\ee^{n\delta+2-y},
 \qquad \delta\ge\eta.
\]
No zero equals \(1\), since \(G(1)=1\).
For \(v=(1-\beta)^{-1}\),
\[
 \Rea v=\frac{1-\Rea\beta}{|1-\beta|^2}\ge\tfrac12,\qquad
 \frac{|v|^2}{\Rea v}=\frac1{1-\Rea\beta}\le\frac1\eta.
\]
Finally \(nm_1=G'(1)/G(1)=d\), so \(m_1\) is real and at least \(1/2\).
\end{proof}

\begin{lemma}[Local analytic estimate]\label{lem:cauchy}
With the normalization above, let \(r>0\) and suppose
\begin{equation}\label{eq:q}
 q=r\sqrt{n+1}\,\exp\bigl(-y+(2n+1)r\bigr)\le\tfrac12.
\end{equation}
Then \(G\) has no zero in \(|w-1|\le r\), and, for every integer \(k\ge2\),
\begin{equation}\label{eq:moments}
 |m_k|\le
 \frac{2\sqrt{n+1}}{n}
       \exp\bigl(-y+(2n+1)r\bigr)r^{1-k}.
\end{equation}
\end{lemma}
\begin{proof}
The disk centered at \(1\) can extend beyond \(\overline\D\);
the needed control there comes from the coefficients.
For \(|u|\le r\),
\[
 |E(1+u)|\le\sqrt{n+1}(1+r)^n\norm E_2
            \le\sqrt{n+1}\,\ee^{-y+nr}.
\]
Solving the ODE from \(G(1)=1\) gives
\[
 G(1+u)=\ee^{\mu u}(1+h(u)),\qquad
 h(u)=\int_0^u\ee^{-\mu s}E(1+s)\,ds.
\]
Using the straight segment and \eqref{eq:mu}, we obtain
\(|h(u)|\le q\le1/2\).
Thus \(G\) is nonzero on that closed disk and on a neighborhood of it.
The analytic function \(F=G'/G-\mu=E/G\) satisfies
\[
 |F(1+u)|\le2\sqrt{n+1}\,\ee^{-y+(2n+1)r}\qquad(|u|\le r).
\]
Cauchy's derivative estimate and logarithmic differentiation give
\[
 \frac{|F^{(k-1)}(1)|}{(k-1)!}
 \le 2\sqrt{n+1}\,\ee^{-y+(2n+1)r}r^{1-k},\qquad
 F^{(k-1)}(1)=(-1)^{k-1}(k-1)!\,n m_k.
\]
This proves \eqref{eq:moments}. There is no restriction \(k\le n+1\):
\(F\) is an analytic rational function on the specified disk.
\end{proof}

\begin{proposition}[Finite certificate]\label{prop:certificate}
Fix \(n\ge1\), \(y,r>0\), and an integer \(K\ge2\).
Define \(\eta\) by \eqref{eq:eta} and set \(a=2\eta\).
Suppose \(0<\eta\le2\), \(\sqrt{n+1}\ee^{-y}\le1\), and
\eqref{eq:q} holds. Put
\begin{equation}\label{eq:B}
 B=\frac{2\sqrt{n+1}\,r}{Kan}
       \exp\left(-y+(2n+1)r+
                       K\log\left(1+\frac ar\right)\right).
\end{equation}
If
\begin{equation}\label{eq:certificate}
 \frac1{K\eta}+B<\frac12,
\end{equation}
then \(s_n\ge\ee^{-y}\).
\end{proposition}
\begin{proof}
Suppose an admissible \(P,\lambda\) had residual at most \(\ee^{-y}\).
For the numbers in Lemma~\ref{lem:normalize},
\[
 |1-av_j|^2
 =1-2a\Rea v_j+a^2|v_j|^2\le1.
\]
Use the bounded Ces\`aro test
\begin{equation}\label{eq:test}
 \Psi(v)=\frac{1-(1-av)^K}{Ka}
       =v+\sum_{k=2}^K
          \frac{(-1)^{k+1}}K\binom Kk a^{k-1}v^k.
\end{equation}
It has zero constant coefficient, linear coefficient one, and
\(|\Psi(v_j)|\le2/(Ka)=1/(K\eta)\).
By Lemma~\ref{lem:cauchy}, its remaining averaged terms have modulus
at most
\begin{align*}
 &\frac{2\sqrt{n+1}\,r}{Kan}\,
       \ee^{-y+(2n+1)r}
       \left[\left(1+\frac ar\right)^K-1-\frac{Ka}{r}\right]
 \le B.
\end{align*}
The triangle inequality and \(m_1\ge1/2\) now imply
\[
 \frac12\le m_1\le
       \left|\frac1n\sum_j\Psi(v_j)\right|+B
       \le\frac1{K\eta}+B,
\]
contrary to \eqref{eq:certificate}.
Each individual residual is therefore larger than \(\ee^{-y}\);
taking its infimum gives the asserted non-strict bound.
\end{proof}

\section{Choice of parameters}
\begin{proof}[Proof of \eqref{eq:asymptotic}]
Fix \(C>C_0\), and choose \(t>0\) so small that
\[
 C^2>\frac43(2+t)\log17.
\]
Set
\[
 y=C\sqrt n,\qquad r=\frac{y}{8n},\qquad
 \eta=\frac{y-\gamma_n}{n},\qquad
 K=\left\lceil\frac{2+t}{\eta}\right\rceil.
\]
For all sufficiently large \(n\), the preliminary conditions of
Proposition~\ref{prop:certificate} hold:
\(\eta\sim C/\sqrt n\), \(q\to0\), and
\(\sqrt{n+1}\ee^{-y}\to0\).
Also \(K\eta\ge2+t\), leaving a fixed positive margin in
\eqref{eq:certificate}. We have
\[
 K\sim\frac{2+t}{C}\sqrt n,\qquad \frac ar\to16,
\]
and hence
\[
 \frac{-y+(2n+1)r+K\log(1+a/r)}{\sqrt n}
 \longrightarrow -\frac{3C}{4}+\frac{2+t}{C}\log17<0.
\]
The prefactor in \eqref{eq:B} is \(O(n^{-1})\).
Thus \(B\to0\), and the finite certificate applies.
\end{proof}

\begin{proof}[Proof of \eqref{eq:explicit}]
Take \(n\ge10^4\) and choose
\[
 y=3\sqrt n,\qquad r=\frac3{8\sqrt n},\qquad
 K=\left\lceil\frac{21}{10\eta}\right\rceil.
\]
First,
\begin{equation}\label{eq:etaexplicit}
 \gamma_n\le\frac3{20}\sqrt n,\qquad
 \frac{19}{20}\frac yn\le\eta\le\frac yn.
\end{equation}
Indeed \(\gamma_n/\sqrt n\) decreases for \(n\ge10^4\), since the
sign of its derivative is that of \(n/(n+1)-\gamma_n<0\).
At \(n=10^4\), \(\gamma_n<10<15\):
\(\ee^8>2^8=256>2\sqrt{10001}\) suffices for the first inequality.
In particular \(0<\eta<2\).
The condition \(\sqrt{n+1}\ee^{-y}<1\) is immediate, and
\[
 q=\frac38\sqrt{1+\frac1n}
       \exp\left(-\frac94\sqrt n+\frac3{8\sqrt n}\right)<\frac12.
\]
Moreover \(K\eta\ge21/10\), so
\[
 \frac12-\frac1{K\eta}\ge\frac1{42}.
\]

From \eqref{eq:etaexplicit} we obtain \(a/r\le16\) and
\(K\le(14/19)\sqrt n+1\). Using \(\log17<3\),
\begin{align*}
 -y+(2n+1)r+K\log(1+a/r)
 &\le\left(-\frac94+\frac{14}{19}\log17\right)\sqrt n
        +\log17+\frac3{8\sqrt n}\\
 &<-\frac3{76}\sqrt n+3+\frac3{8\sqrt n}<0.
\end{align*}
The last expression decreases with \(n\) and is negative at \(10^4\).
For an elementary check of \(\log17<3\), the exponential series
through degree five gives \(\ee^3>18>17\).
Finally \(Ka\ge21/5\), and therefore
\[
 \frac{2\sqrt{n+1}\,r}{Kan}
 \le\frac5{28n}\sqrt{1+\frac1n}<\frac1{5n}.
\]
It follows that \(B<1/(5n)\le1/50000<1/42\).
All conditions of Proposition~\ref{prop:certificate} are satisfied.
\end{proof}

\begin{remark}
The constants \(3\) and \(10^4\) are convenient explicit choices.
The same proof with \(r=sy/n\), \(0<s<1/2\), gives the sufficient
asymptotic condition
\[
 C^2>\frac{2\log(1+2/s)}{1-2s}.
\]
We use \(s=1/8\) to state a simple constant. Neither this expression
nor \(C_0\) is asserted to be the optimal residual constant.
\end{remark}

\section{Exterior power sums and Hardy norms}\label{sec:power}
\begin{corollary}\label{cor:powers}
Let \(|z_j|\ge1\), \(1\le j\le n\), and write
\(M_n=\max_{2\le k\le n+1}|\sum_jz_j^k|\).
Then
\[
 M_n\ge\frac1n\ee^{-3\sqrt n}\qquad(n\ge10^4).
\]
For every \(C>C_0\), uniformly over all such configurations,
\(M_n\ge n^{-1}\ee^{-C\sqrt n}\) for all sufficiently large \(n\).
No distinguished point \(z_1=1\) is needed.
\end{corollary}
\begin{proof}
We use the projected Newton identity of
\cite[Proposition 3.1]{tan}. Put
\(P(t)=\prod_{j=1}^n(1-z_jt)\) and \(p_k=\sum_jz_j^k\).
Then \(P\in\mathcal S_n\), \(P(0)=1\), and \(P'(0)=-p_1\).
Logarithmic differentiation near zero, followed by coefficient
projection onto degrees at most \(n\), gives exactly
\[
 P'-P'(0)P=-\Pi_{\le n}
       \left[\left(\sum_{k=2}^{n+1}p_kt^{k-1}\right)P\right].
\]
Projection is contractive in \(H^2\), and multiplication by a
polynomial \(S\) has \(H^2\) norm at most \(\norm S_\infty\).
Thus
\[
 s_n\le\frac{\norm{P'-P'(0)P}_2}{\norm P_2}
       \le\sum_{k=2}^{n+1}|p_k|\le nM_n.
\]
Apply Theorem~\ref{thm:main}.
\end{proof}

The factor \(1/n\) is retained in these finite statements.
It can be absorbed in a leading square-root exponent with any
strictly larger constant. This corollary strengthens the scale in
\cite{tan}; the earlier qualitative conclusion of \cite{luo} remains
fully credited.

\begin{corollary}\label{cor:hp}
Uniformly for \(1\le p\le\infty\), \(P\in\mathcal S_n\), and
\(\lambda\in\C\),
\[
 \frac{\norm{P'-\lambda P}_{H^p}}{\norm P_{H^p}}
       \ge\frac{\ee^{-3\sqrt n}}{\sqrt{n+1}}\qquad(n\ge10^4).
\]
For every \(C>C_0\), the lower bound \(\ee^{-C\sqrt n}\) holds
for all sufficiently large \(n\), uniformly over the same parameters.
\end{corollary}
\begin{proof}
Write \(R=P'-\lambda P\), of degree at most \(n\).
For \(1\le p\le2\), bounding each Fourier coefficient by
\(\norm R_1\) gives
\[
 \norm R_2\le\sqrt{n+1}\norm R_1
       \le\sqrt{n+1}\norm R_p,\qquad
 \norm P_2\ge\norm P_p.
\]
For \(p\ge2\), use
\[
 \norm R_2\le\norm R_p,\qquad
 \norm P_p\le\norm P_\infty\le\sqrt{n+1}\norm P_2.
\]
In both cases \(\norm R_p/\norm P_p\ge
(\norm R_2/\norm P_2)/\sqrt{n+1}\).
For the asymptotic statement, first apply Theorem~\ref{thm:main}
with a constant strictly between \(C_0\) and \(C\), and absorb
\(\tfrac12\log(n+1)=o(\sqrt n)\).
\end{proof}

\section{Superalgebraic upper bounds}\label{sec:upper}
We prove the upper statement by cancelling a prescribed finite number of
logarithmic-derivative terms at a strict global maximum. All constants in
this section may depend on the prescribed cancellation order \(K\).
They are not asserted to be uniform as \(K\to\infty\).

\begin{lemma}[A finite positive moment cone]\label{lem:cone}
Fix \(K\ge3\). For a rational point \(\zeta=x+iy\), \(x>0\), \(y\ge0\),
define the real column
\[
 V(\zeta)=
 \begin{cases}
  (\zeta^2,\ldots,\zeta^K)^T,&y=0,\\
  (2\Rea\zeta^2,\ldots,2\Rea\zeta^K)^T,&y>0.
 \end{cases}
\]
There is a finite collection of these columns whose cone is
\(\mathbb R^{K-1}\). The collection can contain any prescribed \(K\)
distinct positive rational real points. Its column matrix \(A_\infty\)
admits a strictly positive rational vector \(\alpha^0\) satisfying
\[
 A_\infty\alpha^0=(0,3/2,0,\ldots,0)^T.
\]
\end{lemma}
\begin{proof}
A linear functional nonnegative on every such column would give a
harmonic polynomial \(h(z)=\Rea\sum_{j=2}^K a_jz^j\) nonnegative on
\(\Rea z>0\), by density and conjugation symmetry. If its largest
nonzero degree is \(l\ge2\), choose
\(\theta\in(-\pi/2,\pi/2)\) with \(a_l\cos(l\theta)<0\).
Then \(h(te^{i\theta})<0\) for large \(t\), a contradiction.
Thus the dual cone is zero and the closure of the generated cone is
the whole space. Approximate each positive and negative coordinate
vector by a finite conical combination closely enough that their
convex hull still contains zero in its interior. The finite union of
the columns used in these combinations already generates the whole
space. Append the prescribed real points if necessary.

The resulting rational matrix has full row rank. For each column
\(s_i\), express \(-s_i\) as a nonnegative combination of its columns.
This representation can be rational: a nonempty rational polyhedron
has a rational point. Adding these relations for all \(i\) gives a
rational kernel vector with every coordinate strictly positive.
A rational solution of the displayed full-rank system exists; adding
a sufficiently large positive rational multiple of that kernel vector
makes all of its coordinates positive.
\end{proof}

\subsection{A fixed number of interior zeros}
For a real reciprocal variable \(v>1\), put
\(\beta=1-1/v\) and \(f_v(z)=1-z/\beta\). For \(v=x+iy\),
\(y>0\), use the conjugate-pair factor
\(f_v(z)=(1-z/\beta)(1-z/\overline\beta)\).
Write \(\nu_v=1,2\), respectively, and define
\[
 M_j(v)=\begin{cases}v^j,&v\text{ real},\\2\Rea v^j,&\Im v>0,\end{cases}
 \qquad
 d(v)=\begin{cases}-v^2/(v-1),&v\text{ real},\\
 2\Rea[-v^2/(v-1)],&\Im v>0.\end{cases}
\]
Thus \(d(v)=f_v'(0)-f_v'(1)/f_v(1)\), and
\(\Rea v>1/2\) is equivalent to \(|\beta|<1\).
Use the circle-root baseline
\[
 B(z)=(1+z)^4(1-z+z^2),\qquad
 b_j=\sum_{\beta:\,B(\beta)=0}(1-\beta)^{-j}.
\]
It has degree six, constant coefficient one,
\(B'(0)=B'(1)/B(1)=3\), \(b_2=0\), and \(b_3=-3/2\).

\begin{proof}[Proof of Theorem~\ref{thm:upper}]
Fix \(K\ge3\). Apply Lemma~\ref{lem:cone}, with \(K\) designated
distinct positive real nodes among the \(\zeta_i\). Clear the
denominators of its strictly positive weights by an integer \(Q\),
and put \(m_i=Q\alpha_i^0\). These are positive integer group
multiplicities, fixed independently of \(n\). Set
\[
 M=\sum_i\nu_i m_i,\qquad S=\sum_i m_i M_1(\zeta_i)>0.
\]
The moments of orders \(2,\ldots,K\) are now
\((0,3Q/2,0,\ldots,0)\).
For every large integer \(n\), define
\[
 R=(n/(6Q))^{1/3},\quad \epsilon=R^{-1},\quad
 c_n\equiv n-M\pmod6,\quad |c_n-2SR|\le3,
 \qquad b_n=(n-M-c_n)/6.
\]
Choose a nearest integer in the specified congruence class; both
\(b_n,c_n\) are positive eventually. Then
\(t_n=c_n/R=2S+O(\epsilon)\) and
\(u_n=b_n/R^3=Q+O(\epsilon^2)\).

Let \(\xi_i\) vary near \(\zeta_i\), varying only the \(K\)
designated real nodes. Define
\(d_\epsilon(\xi)=-\xi^2/(\xi-\epsilon)\) for a real node,
and twice its real part for a pair. Consider the \(K\) equations
\begin{align}
 0&=t/2+\sum_i m_i d_\epsilon(\xi_i),\label{eq:defect-slope}\\
 0&=t\epsilon/4+\sum_i m_i M_2(\xi_i),\label{eq:defect-two}\\
 0&=b_j u\epsilon^{j-3}+t2^{-j}\epsilon^{j-1}
       +\sum_i m_i M_j(\xi_i),\qquad 3\le j\le K.
       \label{eq:defect-high}
\end{align}
They vanish at \((\epsilon,t,u,\xi)=(0,2S,Q,\zeta)\).
The Jacobian in the designated real positions has columns
\[
 m_i(-1,2\zeta_i,3\zeta_i^2,\ldots,K\zeta_i^{K-1})^T.
\]
It is a nonsingular scaled Vandermonde matrix. The real
implicit-function theorem at this fixed limiting point gives
\(\xi_{i,n}=\zeta_i+O(\epsilon)\) solving the equations at
\((\epsilon,t_n,u_n)\). Put \(v_{i,n}=R\xi_{i,n}\) and
\[
 P_n(z)=B(z)^{b_n}(1+z)^{c_n}\prod_i f_{v_{i,n}}(z)^{m_i}.
\]
For large \(n\), its degree is exactly \(n\), its constant
coefficient is one, and its coefficients are real. Exactly \(M\)
zeros, counted with multiplicity, lie strictly inside the disk;
the other \(n-M\) lie on the circle. The interior zeros satisfy
\(|1-\beta_{i,n}|\asymp1-|\beta_{i,n}|\asymp R^{-1}\).
Multiplying \eqref{eq:defect-slope}--\eqref{eq:defect-high} by
\(R,R^2,\ldots,R^K\), respectively, gives the exact identities
\[
 \lambda_n:=P_n'(0)=P_n'(1)/P_n(1),\qquad
 \sum_{\beta:\,P_n(\beta)=0}(1-\beta)^{-j}=0
       \quad(2\le j\le K).
\]
In particular integer rounding has been corrected, and
\(\lambda_n=3b_n+c_n/2+R\sum_i m_iM_1(\xi_{i,n})=O(n)\).

We give global estimates before using the cancellations locally.
For \(s=\sin^2(\theta/2)\),
\[
 |B(\ee^{i\theta})/B(1)|^2
 =\cos^6(\theta/2)\cos^2(3\theta/2)\le\ee^{-3s},
 \qquad |(1+\ee^{i\theta})/2|^2\le\ee^{-s}.
\]
For a real group \(v=x>1\), its normalized squared modulus is
\(1+4x(x-1)s\le\ee^{T(v)s}\), with \(T(v)=4x(x-1)\).
For a pair set \(Y=y^2\), \(U=x^2+Y\), \(V=(x-1)^2+Y\).
Its normalized squared modulus is
\[
 1+8(x^2-x-Y)s+16UVs^2
 \le(1+4\sqrt{UV}s)^2\le\ee^{4(U+V)s}.
\]
Indeed \(UV-(x^2-x-Y)^2=(2x-1)^2Y\ge0\).
Set \(T(v)=4(U+V)\) for this group. The fixed number of
interior factors has total cost \(\sum_i m_iT(v_{i,n})=O(R^2)\).
Since \(R\asymp n^{1/3}\),
\[
 3b_n+c_n-\sum_i m_iT(v_{i,n})=n/2-O(R^2)\ge\kappa n
\]
for some fixed \(\kappa>0\) and all sufficiently large \(n\).
It follows globally that
\begin{equation}\label{eq:defect-global}
 |P_n(\ee^{i\theta})/P_n(1)|^2
 \le\exp[-\kappa n\sin^2(\theta/2)].
\end{equation}

The zero-free neighborhood now shrinks with \(n\). Choose a fixed
\(\delta>0\) so that on \(|z-1|\le\delta/R\) all
\(|v_{i,n}(z-1)|\le1/2\), and the baseline zeros are excluded.
Expanding \(v/(1+v(z-1))\), the exact cancellations above give
\begin{equation}\label{eq:defect-local}
 |P_n'(z)/P_n(z)-\lambda_n|
 \le C_K(n+R^{K+1})|z-1|^K
 \le C_KR^{K+1}|z-1|^K.
\end{equation}
The geometric-series remainder is uniform on this shrinking disk.
There are \(O(n)\) baseline reciprocal variables of bounded size
and only \(M\) interior variables of size \(O(R)\).
By real symmetry the first angular derivative of the normalized
log modulus vanishes at zero; its second derivative on a suitably
smaller arc \(|\theta|\le\delta'/R\) is \(O(n+R^2)=O(n)\).
Hence the normalized squared denominator is at least \(c n^{-1/2}\),
by integration over \(|\theta|\le n^{-1/2}\).
Combining \eqref{eq:defect-global} and \eqref{eq:defect-local},
the local squared numerator is at most
\[
 C_KR^{2K+2}\int_{\mathbb R}\theta^{2K}\ee^{-c n\theta^2}\,d\theta
 \le C_K R^{2K+2}n^{-K-1/2}.
\]

Outside the small arc use the polynomial product derivative.
Normalized derivatives of a baseline or extra factor are bounded;
those of a real interior factor are \(O(R)\), and those of a pair
are \(O(R^2)\). Removing a baseline factor reduces the positive
global cost by three, removing an extra factor reduces it by one,
and removing an interior group increases it by its positive \(T(v)\).
Thus every derivative term is bounded by a polynomial in \(n\)
times the same Gaussian bound, up to a constant. On the outer arc
\(n\sin^2(\theta/2)\ge c n/R^2\asymp n^{1/3}\).
Its squared integral is therefore at most
\(\operatorname{poly}(n)\ee^{-c n^{1/3}}\), negligible compared
with every fixed inverse power. This estimate uses no rational
quotient at a boundary zero. Dividing by the denominator proves
\[
 \rho(P_n)^2\le C_K R^{2K+2}n^{-K}
 \le C_K n^{(2-K)/3},\qquad
 \rho(P_n)\le C_K n^{(2-K)/6}.
\]
Choose \(K\ge6A+2\). Then \(M_A=M\) is independent of \(n\),
as required. The proof does not give effective bounds on \(M_A\),
\(C_A\), or the starting degree as functions of \(A\).
\end{proof}

\begin{theorem}[One interior zero suffices]\label{thm:one}
For every sufficiently large integer \(n\), there is a real
degree-\(n\) polynomial \(P_n\), with \(P_n(0)=1\), exactly one
zero strictly inside the unit disk and all remaining zeros on the
circle, such that \(\rho(P_n)=O(n^{-1/6})\).
The interior zero is real and positive, with distance and radial gap
from 1 of order \(n^{-1/3}\). In conjunction with
Theorem~\ref{thm:circle}, the minimum number of interior zeros
permitting a sequence with vanishing residual is exactly one.
No optimal rate for one interior zero is asserted.
\end{theorem}
\begin{proof}
We use two adjustable circle angles in the degree-34 baseline
\[
 B_{a,c}(z)=(1-2az+z^2)^{15}(1-2cz+z^2)^2,
 \qquad -1<a,c<1.
\]
Its slope difference and second reciprocal moment are, respectively,
\[
 D(a,c)=-30a-4c-17,\qquad
 T(a,c)=-\frac{15a}{1-a}-\frac{2c}{1-c}.
\]
Both vanish at \((a_*,c_*)=(-2/3,3/4)\). Write
\(n-1=34m+k\), \(0\le k\le33\), and let \(v=v_n>3/2\)
be the unique solution of
\begin{equation}\label{eq:single-cubic}
 v^3-\tfrac32v^2=(n-1)/4.
\end{equation}
The left side is strictly increasing on this interval, and
\(v\sim(n/4)^{1/3}\). Set \(\beta=1-1/v\in(0,1)\).
For large \(n\), \(m>0\). Define the real targets
\[
 d_n=\frac{v^2/(v-1)-k/2}{m}=O(n^{-2/3}),\qquad
 t_n=\frac{-v^2-k/4}{m}=O(n^{-1/3}).
\]
The following explicit formulas select circle parameters solving
\(D(a_n,c_n)=d_n\), \(T(a_n,c_n)=t_n\):
\begin{align*}
 L&=(51+d_n)/2,\qquad U=17-t_n,\\
 X&=\frac{UL+221+\sqrt{(UL-289)(UL-169)}}{30U},\\
 Y&=(L-15X)/2,\qquad a_n=1-X,\qquad c_n=1-Y.
\end{align*}
Indeed these equations are equivalent to
\(15X+2Y=L\), \(15/X+2/Y=U\); elimination gives
\(15UX^2-(UL+221)X+15L=0\). The positive-square-root branch
above tends to \(X=5/3,Y=1/4\), since
\(UL\to867/2\) and the square root tends to \(391/2\).
The denominators and discriminant stay positive, and
\(a_n,c_n\in(-1,1)\) eventually, with deviations
\(O(n^{-1/3})\) from \((a_*,c_*)\). These statements are
uniform over the 34 values of \(k\).

Put
\[
 P_n(z)=B_{a_n,c_n}(z)^m(1+z)^k(1-z/\beta).
\]
It has exactly the claimed degree, constant term and zero locations.
The two target equations give the exact identities
\[
 P_n'(0)=P_n'(1)/P_n(1)=:\lambda_n,
 \qquad \sum_{j=1}^n v_j^2=0,
 \qquad v_j=(1-\beta_j)^{-1}.
\]
Here the extra roots at \(-1\) contribute \(k/2\) to the slope
difference and \(k/4\) to the second moment; the interior root
contributes \(-v^2/(v-1)\) and \(v^2\), respectively.
For a conjugation-symmetric multiset of \(N\) circle roots
avoiding 1, its reciprocal moments satisfy
\begin{equation}\label{eq:circle-third}
 M_3^{\circ}=\tfrac32 M_2^{\circ}-N/4.
\end{equation}
This follows for each pair \(1/2\pm iy\), and also for the real
reciprocal \(1/2\). Conjugation symmetry is essential and holds
here. Taking \(N=n-1\) and \(M_2^{\circ}=-v^2\),
\eqref{eq:single-cubic} gives \(\sum_jv_j^3=0\) exactly.
Moreover \(\lambda_n=(n-1)/2+v=O(n)\).

We verify the global norm control for this new baseline. In
\(s=\sin^2(\theta/2)\), its limiting normalized squared modulus is
\[
 H(s)=(1-6s/5)^{30}(1-8s)^4,\qquad 0\le s\le1.
\]
Apart from its zeros \(1/8,5/6\), the only interior stationary
point is \(5/24\). Its values there and at 1 are
\[
 H(5/24)=(3/4)^{30}(2/3)^4<1,\qquad H(1)=7^4/5^{30}<1.
\]
Since \(H(0)=1\), zero is the unique global maximum.
Also \(\log H(s)=-68s+O(s^2)\). The nearby circle parameters
retain uniform negative curvature near zero and, by compactness,
a strict modulus gap away from zero. Therefore, for some fixed
\(\kappa>0\) and all sufficiently large \(n\),
\[
 |B_{a_n,c_n}(\ee^{i\theta})/B_{a_n,c_n}(1)|^2
 \le \ee^{-\kappa s}\qquad(-\pi\le\theta\le\pi).
\]
The interior factor contributes
\(1+4v(v-1)s\le\exp(4v(v-1)s)\), whereas the \(-1\)
factors contribute at most \(\ee^{-ks}\). Since
\(m\sim n/34\) and \(v^2=O(n^{2/3})\),
\(\kappa m+k-4v(v-1)\ge c n\) eventually. Thus
\begin{equation}\label{eq:single-global}
 |P_n(\ee^{i\theta})/P_n(1)|^2
 \le\exp[-c n\sin^2(\theta/2)]
\end{equation}
globally, for a fixed \(c>0\).

The circle reciprocals are uniformly bounded and only the interior
reciprocal is large. On \(|z-1|\le\delta/v\), for a fixed small
\(\delta>0\), all denominators below stay separated from zero.
The exact slope and moment cancellations imply
\[
 F_n(z):=P_n'(z)/P_n(z)-\lambda_n
 =-(z-1)^3\sum_{j=1}^n\frac{v_j^4}{1+v_j(z-1)}.
\]
Consequently \(|F_n(z)|\le C(n+v^4)|z-1|^3
\le Cn^{4/3}|z-1|^3\). On a smaller angular window of radius
\(\delta'/v\), real symmetry makes the first derivative of the
normalized log modulus zero at zero, and its second derivative is
\(O(n+v^2)=O(n)\). As \(n^{-1/2}=o(v^{-1})\), integration
on \(|\theta|\le n^{-1/2}\) bounds the normalized squared
denominator below by \(c'n^{-1/2}\). The local squared numerator
is at most
\[
 Cn^{8/3}\int_{\mathbb R}\theta^6\ee^{-c'n\theta^2}\,d\theta
 =O(n^{-5/6}).
\]
On the outer arc differentiate the polynomial product, avoiding
logarithmic-derivative poles. The normalized derivative of a baseline
factor is uniformly bounded, its multiplicity is \(O(n)\), and
the normalized interior derivative is \(v\). Removing a baseline
loses only \(\kappa\) in the global exponent; removing the interior
factor improves it. The extra-factor multiplicity is bounded.
The term \(-\lambda_nP_n\) costs at most another \(O(n)\).
Thus the outer squared integral is bounded by a polynomial in
\(n\) times \(\exp(-c'n/v^2)=\exp(-c''n^{1/3})\).
Dividing both numerator estimates by the same denominator gives
\(\rho(P_n)^2=O(n^{-1/3})\), as required.

Finally, Theorem~\ref{thm:circle} excludes vanishing residuals with
zero interior roots; this construction permits them with one.
This proves the zero-count threshold, not a sharp rate or the
superalgebraic conclusion with one root. No numerical starting
degree or norm constant is claimed.
\end{proof}


\section{A fixed radial gap forces a polynomial barrier}\label{sec:radial}
The upper construction uses the full closed disk. A fixed strict
root radius gives a complementary obstruction, by the same local
estimate and bounded test as the lower-bound proof.

\begin{proposition}\label{prop:radial}
For \(0\le q<1\), let \(s_n(q)\) denote the infimum defining \(s_n\)
with all zeros restricted to \(|z|\le q\). Set
\[
 a=1-q^2,\quad K=\lceil4/a\rceil,\quad
 \Delta=\frac1{1+q}-\frac{1+q^K}{Ka}>0,\qquad
 c_q=\min\left\{\frac1{\sqrt2},
       \frac{\Delta Ka}{2\sqrt2\ee^3(1+a)^K}\right\}.
\]
Then \(s_n(q)\ge c_q n^{3/2-K}\) for every \(n\ge41\).
\end{proposition}
\begin{proof}
Write \(\rho\) for an admissible residual. If
\(\rho>1/\sqrt{n+1}\), the conclusion is immediate.
Otherwise the maximum normalization gives \(|\mu|\le n+1\).
Take \(r=1/n\). The smallness parameter in the proof of
Lemma~\ref{lem:cauchy} is at most \(\ee^3/n<1/2\).
That local proof does not need the separate zero-separation
parameter \(\eta\), and gives
\[
 |m_j|\le\frac{2\sqrt{n+1}}n\rho\ee^3 r^{1-j}\qquad(j\ge2).
\]
The exact image of \(|\beta|\le q\) under \(v=(1-\beta)^{-1}\)
is \(|v-1/a|\le q/a\). Thus \(|1-av|\le q\) and
\(\Rea v\ge1/(1+q)\).
Use \(\Psi(v)=[1-(1-av)^K]/(Ka)\), whose modulus is at most
\((1+q^K)/(Ka)\). Estimating its higher averaged terms gives
\[
 \Delta\le\frac{2\ee^3\sqrt{n+1}\rho}{Kan^2}
       [(1+an)^K-1-Kan]
 \le\frac{2\ee^3\sqrt{n+1}\rho}{Kan^2}(1+an)^K.
\]
Now \(\sqrt{n+1}\le\sqrt{2n}\) and \(1+an\le(1+a)n\).
The resulting lower bound is the stated one. Finally \(Ka\ge4\),
\(q^K<1\), and \(1/(1+q)>1/2\) ensure \(\Delta>0\).
\end{proof}

In particular, any single sequence of admissible degree-\(n\)
polynomials with residual \(o(n^{-A})\) for every fixed \(A>0\)
must have its maximum root modulus tending to one. Otherwise an
infinite subsequence would satisfy a fixed polynomial lower bound.
No sharp radial exponent is asserted.

\section{Scope and verification}
The uniform slope formulation bounds approximate eigenvectors of
polynomial differentiation under a zero-location constraint.
It does not describe the full unconstrained pseudospectrum.
For example, for \(u=n/4\) and
\(T_n(z)=\sum_{k=0}^n(uz)^k/k!\),
\[
 T_n'-uT_n=-\frac{u^{n+1}}{n!}z^n,\qquad
 \frac{\norm{T_n'-uT_n}_2}{\norm{T_n}_2}
       \le\frac n4\left(\frac{\ee}{4}\right)^n,
\]
using \(\norm{T_n}_2\ge1\) and \(n!\ge(n/\ee)^n\).
Thus a root constraint is essential for a bound of our form.

The lower and upper disk bounds do not match in order.
Whether the restricted exterior power-sum minima stay bounded
away from zero remains unresolved here, as does the optimal order
of \(s_n\) and \(r_n\).
The superalgebraic upper bound is not a stretched-exponential upper
bound, and its constants are not optimized. The minimal interior-root
count permitting collapse is one, but the optimal rate with one such
root is not established. We do not prove superalgebraic decay with
one interior root. The elementary radial
barrier is a consequence of the same local method.
We do not assert a theorem for arbitrary starting exponents,
weighted power sums, or \(H^p\) with \(p<1\).
The elementary norm comparison in Corollary~\ref{cor:hp} is a
consequence of the main estimate.

The analytic proof of the explicit threshold uses only the displayed
elementary inequalities. It does not rely on optimization or sampling.
The supplementary diagnostics check exact projected Newton identities
over Gaussian rationals, the bounded test on admissible reciprocal
zeros, binomial identities, and numerical instances of the finite
certificate. Rational checks of constants and high-precision
evaluations are included for reproducibility. The circle-root family
is checked with complete exact coefficient Parseval ratios. The
one-interior-root family has symbolic and high-precision checks of
the cubic, slope, moments and circle parameters, without quadrature.
The independent audit covers all 34 degree classes at two scales.
These finite checks do not certify the full norm asymptotics.
Symbolic identities support the constructions. Finite moment
searches motivated the general construction but are not premises
of its arbitrary-order proof.
These checks do not establish historical priority or replace the proof.

The comparison uses the exact primary-source versions listed below.
A targeted search did not identify the bounds in this exact setting;
exhaustive priority certification is not claimed. The accessible
Komarov versions listed below were used for the stated comparison.
A related 2020 journal text was not available in full in this review;
no exclusion of every possible bibliographic overlap is claimed.
AI tools assisted exploration, drafting and computational checking.
Model review is distinct from independent human peer review and
formal proof verification.

\begin{thebibliography}{9}
\bibitem{tan}
X. Tan, Q. Wang, W. Huang and K. Chen,
\emph{Residual bounds for Schur-stable polynomials},
arXiv:2608.02043v3, 8 August 2026.
\url{https://arxiv.org/abs/2608.02043v3}.
\bibitem{luo}
Y. Luo, R. Yang and K. Zhu,
\emph{Exterior power sums},
arXiv:2607.22017v1, 24 July 2026.
\url{https://arxiv.org/abs/2607.22017v1}.
\bibitem{komarovrate}
M. A. Komarov,
\emph{Rate of approximation by logarithmic derivatives of polynomials
whose zeros lie on a circle}, arXiv:1804.08918v1, 24 April 2018.
\url{https://arxiv.org/abs/1804.08918v1}.
\bibitem{komarovmoments}
M. A. Komarov,
\emph{Caratheodory type representation with unit weights and related
approximation problems}, arXiv:1807.06499v1, 17 July 2018 (Russian).
\url{https://arxiv.org/abs/1807.06499v1}.
\bibitem{erdelyi}
T. Erd\'elyi, D. P. Hardin and E. B. Saff,
\emph{Inverse Bernstein inequalities and min-max-min problems on
the unit circle}, arXiv:1307.4056v1, 15 July 2013.
\url{https://arxiv.org/abs/1307.4056v1}.
\bibitem{busch}
P. Busch, J. Kiukas and R. F. Werner,
\emph{Sharp uncertainty relations for number and angle},
arXiv:1604.00566v2, \S VI.A.2.
\url{https://arxiv.org/abs/1604.00566v2}.
\bibitem{dlmf}
NIST Digital Library of Mathematical Functions,
\S\S28.2 and 28.8, especially Eq.~(28.8.1).
\url{https://dlmf.nist.gov/28.2}; \url{https://dlmf.nist.gov/28.8}.
Accessed 18 October 2026.
\end{thebibliography}
\end{document}
